\scanpage{131}
Brives-Charensac, le 08/04/1986\\
Monsieur le Préfet,\\

Avec tous ceux qui oeuvrent pour une prise en charge personnelle de la santé, que cela soit celle des humains ou celle des animaux, les deux étant intimement liées (il est bien évident que les lois veillant à la protection du cheptel ont en vue la santé publique), les homœopathes se sentent fortement concernés par le procès qui est fait à Monsieur COULARDEAU pour refus de vacciner des animaux traités par homœopathie :

- Ce procès revient à nier purement et simplement, de la part des pouvoirs publics, l'homœopathie, sa valeur préventive, donc sa valeur tout \unsure{\uline{court}|court}, puis qu'elle se donne essentiellement comme une médecine de terrain.

- Il choque la conscience de tous ceux qui ont fait l'expérience du retour à une immunité naturelle solide après traitement homœopathique, et veulent en faire bénéficier leurs enfants.

- Il insulte les médecins qui se donnent le mal, pour le bien de leur prochain, d'acquérir cette difficile technique.

- Il choque tous les amis de la liberté de pensée, renouvelant une sorte d'inquisition qui cherche à obtenir de force l'adhésion de la conscience à des dogmes indémontrables.

Des événements récents publiés par la presse confirment que l'on ne touche pas impunément à l'appareil immunitaire d'un organisme vivant, et si les laboratoires scientifiques ne peuvent pas prouver (ou ne veulent pas) qu'une agression vaccinale puisse faire basculer un organisme en difficulté transitoire, ils ne peuvent pas non plus nous démontrer le contraire.

Doit-on, en république française, supporter de voir Mr COULARDEAU payer une amende pour refus d'appartenir à la religion d'état ?

Si la loi impose à chaque éleveur de maintenir les animaux en bonne santé, il semble prouvé que Mr COULARDEAU a satisfait à la loi plus que nombre d'éleveurs qui vaccinent leurs bêtes. Il a donc satisfait à l'esprit de la loi plus que ceux qui satisfont à la lettre.

\hfill .../...

% First copy of the closing, at left when the signature sheets are right side up.
En attendant que soit réformée la lettre au nom de l'esprit qui l'a inspirée, nous vous prions, Monsieur le Préfet, de bien vouloir user de votre autorité pour que Mr COULARDEAU soit jugé en fonction de cet esprit.

Dans cette confiance, nous vous prions de croire, Monsieur le Préfet, à l'expression de notre haute considération.\\

\handwritten{\unsure{A.|J.} Fayeton}\\

\handwritten{D\textsuperscript{r} Alain \unsure{Reynaud|Reynauld} \qquad \unsure{R}\\
(Montpellier)}\\

\handwritten{D \unsure{Diendomné|Dieudonné} \qquad \illegible{1}\\
le PUY.}\\

\handwritten{D\textsuperscript{r} Denis \unsure{GRANGEOUT|GRANGEON} \unsure{13} cours Victor HUGO 43000 LE PUY \qquad \unsure{Grangeon}}\\

\handwritten{Dr Damien DU PERRON \unsure{Ch|ch} \unsure{fr|St} Jean. St \unsure{Todard|Jodard}. \unsure{42590|42390} \unsure{NEULISE} \qquad \unsure{Du Perron}}\\

\handwritten{D\textsuperscript{r} \unsure{Naouri|Naoury} 60 Rue du P\textsuperscript{r} E Herriot Lyon \unsure{69002} \qquad \illegible{1}}\\

\handwritten{Dr G. \unsure{LOUTAN|COUTAN} 7. \unsure{Conseil-Général} 1205 GENÈVE \qquad \illegible{1}}\\

\handwritten{Dr P. \unsure{Harzé|Harzer} ch Ste \unsure{d'Hudard|d'Hudson} 42550 \unsure{Nervisé|Nervieux} \qquad \illegible{1}}\\

\handwritten{Dr. F. Vidal 34 rue de l'Horloge. 63200 Riom. \qquad \unsure{F V}}\\

\handwritten{Dr \unsure{Ch|de} Louise \unsure{Antony|Autton}, 26, rue de la \unsure{Griselle|Gabelle} 1228 \unsure{Genève|Genève-} Genève}\\

% Second copy of the closing, with a different set of signatures.
\mspage{102}
En attendant que soit réformée la lettre au nom de l'esprit qui l'a inspirée, nous vous prions, Monsieur le Préfet, de bien vouloir user de votre autorité pour que Mr COULARDEAU soit jugé en fonction de cet esprit.

Dans cette confiance, nous vous prions de croire, Monsieur le Préfet, à l'expression de notre haute considération.\\

\handwritten{D\textsuperscript{r} FAYETON. BRIVES CHARENSAC. 43700. \qquad \unsure{A.|J.} Fayeton.}\\

\handwritten{D\textsuperscript{r} Marc \unsure{DEGRAIX|DEGRAUX} de LAVISON \qquad \unsure{Alfred|Dfrid|David}\\
11 rue Luizet \qquad \unsure{d'Hillery|d'Hilley}\\
69130 ECULLY.\\
\unsure{13660} Orgon}\\

\handwritten{D\textsuperscript{r} LAURENT.A. \qquad \unsure{A Laurent}\\
13 Av. des 4 Chemins\\
92330 Sceaux}\\

\handwritten{DR \unsure{PETITJEAN}. Lyon \qquad \unsure{Petitjean}}\\

\handwritten{\unsure{SAUSSUS|SAUSSIS} \quad SEYSSINS 38170 \qquad \unsure{Saussus}}\\

\handwritten{Dr VIGNERON J.L CONDRIEU \qquad \unsure{Vigneron}}\\

\handwritten{D\textsuperscript{r} \unsure{Annie|Anne} \unsure{Jacquin|Jacquet} \unsure{Barre} \qquad \unsure{A J Barre}\\
Les \unsure{Sabines|Sabins} 69130 Ecully}\\

\handwritten{\unsure{Didier} \unsure{Cottencin|Cottereau} \qquad \illegible{1}\\
57 \unsure{bd|fbg} St Michel - 75005 Paris}\\

\handwritten{D\textsuperscript{r} \uline{\unsure{M. de K.|M. de K. S}}\\
32 rue St Joseph 31400\\
Toulouse}\\

\handwritten{D\textsuperscript{r} \unsure{MARZIALE|MARZIAKE} \qquad \illegible{1}\\
365 rue de Vaugirard\\
75015 Paris}\\

\handwritten{D\textsuperscript{r} \unsure{BARRAT|BARRAS} \qquad \illegible{1}\\
3 quai \unsure{Julliette|Juliette|Gailleton}}\\

\handwritten{Dr \unsure{MEUNIER|MECHAIN} 13 Av des 4 chemins \qquad \unsure{Médian|Médiau}\\
92330 SCEAUX}\\

\handwritten{\unsure{Dr} \unsure{Vivier|Viver} \qquad '' \qquad \illegible{1}}\\

\handwritten{\deleted{\unsure{Dr}} Roger \unsure{Vidal|Widal}. La Rochelle \qquad Roger \unsure{Widal|Vidal}.}\\

\handwritten{Dr \unsure{Vully|Vally}. 69002. \qquad \illegible{1}}\\
% archive sheet number: 131
